\section{Model I: Physical Environment and Cooling}
 
\subsection{Model Description}
 
Model I follows heat, moisture, airflow, electricity, and water. It starts with outdoor temperature, humidity, pressure, and ventilation. It then calculates the outlet air from each cooling device and mixes that air into the waiting area. The Fan System changes the airspeed felt by a passenger but does not directly lower the zone air temperature; its modeled benefit comes from increased air movement at the person~\cite{cbe_fans_2023}. The Evaporative Cooler and Large Evaporative Cooler lower the air temperature, and the zone model determines how much of that outlet cooling remains after mixing with outdoor air.
 
\subsection{Primary Equations}
 
The first step converts relative humidity into a humidity ratio $w$, which is the mass of water vapor per unit mass of dry air, using the ASHRAE relationships implemented in PsychroLib~\cite{psychrolib_volume}:
\begin{equation}
w = \frac{0.621945\, p_v}{p - p_v}, \qquad p_v = \frac{RH}{100}\, p_{sat}(T)
\end{equation}
The saturation pressure $p_{sat}$ is calculated from the ASHRAE correlation. The wet-bulb temperature $T_{wb}$ is found by solving the psychrometric energy balance. This matters because direct evaporative cooling cannot cool air below its entering-air wet-bulb temperature in this ideal direct-evaporation model~\cite{pnnl_evaporative}.
 
For the selected saturation-limited model, the outlet temperature drop is
\begin{equation}
\Delta T = \min\left[\Delta T_{rated},\; T_{out} - T_{wb}\right], \qquad T_{supply} = T_{out} - \Delta T
\end{equation}
The cooler is treated as an approximately constant-enthalpy process, where the enthalpy of moist air in kJ/kg dry air is\cite{psychrolib_volume}
\begin{equation}
h = 1.006\,T + w\,(2501 + 1.86\,T)
\end{equation}
The supply humidity ratio is chosen so that the inlet and outlet enthalpies match. If that result would exceed 100\% relative humidity, the model moves the outlet back to the saturation boundary. The mass of evaporated water is
\begin{equation}
m_{water} = \dot{m}_{dry\ air}\,(w_{supply} - w_{out})\,(3600\,\Delta t)
\end{equation}
The waiting area is solved with steady mass, heat, and moisture balances. In simple terms, the air and moisture flowing into the zone must equal what flows out. For temperature, each airflow carries approximately $\dot{m}c_pT$. Electrical energy is
\begin{equation}
E = P\,\Delta t
\end{equation}
 

\subsection{Data Sources and Reference Cases}
The main weather comparison uses simultaneous NASA POWER hourly temperature, relative humidity, and pressure for June--August 2025 at the three city reference coordinates~\cite{nasa_hourly_summer_2025}. These are gridded reference conditions, not measurements at a shelter. For each city, the study selects the hottest hour, then the most humid and driest distinct hours within the hottest 10\% of records, using humidity ratio; this selection rule is our scenario design, not a NASA classification.

The transit references are Rosa Parks Plaza and West End Station in Dallas; the Grove Hall clear-roof control shelter and Mattapan outdoor trolley waiting area in Boston; and the bus and light-rail waiting areas at Central Avenue/Camelback in Phoenix~\cite{dart_rosa,dart_west,boston_shade,bwa_mattapan,valley_bus,valley_rail}. The larger-area references are Dallas College Mountain View Student Courtyard and the Love Field curbside canopy, Boston University's George Sherman Union Courtyard and Logan Terminal C canopy, and ASU's Delph Courtyard and Sky Harbor Terminal 3 canopy~\cite{dallas_college_student_courtyard,dal_lovefield_canopy,bu_gsu_courtyard,bos_logan_canopy,asu_delph_courtyard,phx_skyharbor_canopy}. These sources establish location identities, not measured airflow or surveyed geometry.

Ventilation and baseline passenger airspeed are scenario estimates. The wind-driven relation $Q=C_w A U$ and the urban wind-profile exponent are drawn from EnergyPlus; equivalent apertures and passenger sheltering factors are analyst assumptions~\cite{energyplus_wind_ventilation,energyplus_242_surface}. The estimates use the city summer-mean wind references, and the enlarged courtyard/canopy openings are documented project choices~\cite{nasa_2026_dallas_summer,nasa_2026_boston_summer,nasa_2026_phoenix_summer,airflow_scenario_2026,case_study_airflow}.

Contest systems A, B and C use powers of 0.25, 1.2 and 2.5 kW, with stated maxima of 3$^\circ$C perceived cooling, 5$^\circ$C air cooling and 8$^\circ$C air cooling, respectively~\cite{comap_2026_problem_b}. Reference products are the 20-inch AirEye fan and Portacool APEX 2000/4000 coolers. The recorded reference prices are \$1,499, \$4,159 and \$4,939; the selected cooler flows are 9,000 and 13,350 CFM~\cite{aireye_20_price,portacool_2000,portacool_4000}. The fan exposure uses a selected 357 ft/min contour value from the AirEye sheet~\cite{aireye_2021_specs}. These reference specifications do not establish the contest devices' actual site performance or a measured passenger-position fan speed.

\subsection{Key Assumptions and Justifications}
 
\begin{enumerate}
  \item \textbf{The waiting area is one well-mixed zone.} Temperature and humidity are treated as uniform. This is reasonable for a first comparison because exact spatial airflow maps are unavailable. It is not a claim that every seat receives the same cooling.
  \item \textbf{The default model includes mixing but omits surfaces and internal heat gains.} Unknown roof temperatures, wall heat transfer, solar gain, lighting, and passenger heat are not set to invented values. The simplification keeps the calculation reproducible while making it less site-specific.
  \item \textbf{Evaporative cooling follows a constant-enthalpy process.} This is a standard ideal model for direct evaporation. It provides a physical upper limit, but real units may perform below it.
  \item \textbf{Cooling is limited by the wet-bulb temperature and saturation.} This prevents impossible outlet states. It is justified by thermodynamics and is more credible than applying the rated temperature drop in all weather.
  \item \textbf{Rated power and airflow represent the candidate devices.} These values allow equal comparisons, but actual power, throw distance, maintenance condition, and placement may differ.
  \item \textbf{Each modeled period lasts one hour.} This creates comparable energy and water totals. It does not represent a full operating season.
\end{enumerate}
 